\begin{remark}[主特征支的椭圆一致化：\texorpdfstring{$\lambda_+(y)$}{lambda+(y)} 的 Weierstrass \texorpdfstring{$\wp$}{wp} 表示]\label{rem:xi-terminal-zm-lambda-plus-wp-uniformization}
在结论 \ref{con:xi-terminal-zm-elliptic-normalization} 的椭圆模型
\[
E:\ Y^2=X^3-X+\frac14
\]
上，复解析一致化给出存在周期格 \(\Lambda\subset\CC\) 与 Weierstrass 函数 \(\wp_\Lambda\) 使得
\[
E(\CC)\cong \CC/\Lambda,\qquad
X=\wp_\Lambda(z),\qquad
Y=\frac{\wp_\Lambda'(z)}{2},
\]
并满足微分方程
\[
\bigl(\wp_\Lambda'(z)\bigr)^2=4\wp_\Lambda(z)^3-4\wp_\Lambda(z)+1
\qquad(g_2=4,\ g_3=-1).
\]
在该一致化下，物理主支重量函数（结论 \ref{con:xi-terminal-zm-y-quadratic-extension-minpoly}）
\[
y=X^2+Y-\frac12
\]
提升为椭圆函数
\[
y(z)=\wp_\Lambda(z)^2+\frac{\wp_\Lambda'(z)-1}{2}.
\]
令 \(\lambda(z):=X(z)=\wp_\Lambda(z)\)。
由结论 \ref{con:xi-terminal-zm-elliptic-normalization} 的双有理互逆 \(\nu:E\to\mathcal C\)，
对任意 \(z\in\CC/\Lambda\) 有 \((\lambda(z),y(z))\in\mathcal C(\CC)\)，即
\[
\Pi\bigl(\lambda(z),y(z)\bigr)=0.
\]
反之，\(\Pi(\lambda,y)=0\) 的正规化在复解析意义下同构于 \(E(\CC)\cong\CC/\Lambda\)，
因此其任一解析支均可由上述椭圆函数参数化统一给出。
因此 \(\pi_y\) 的有限分歧点在一致化参数 \(z\) 下等价于临界点条件 \(\frac{\dd}{\dd z}y(z)=0\)；
而 \(y=\infty\) 对应于 \(z\to 0\) 的四阶极点结构，与 \((y)_\infty=4O\) 的全分歧极点描述一致。
令 \(P_1=(2,-\tfrac52)\in E(\QQ)\)（满足 \(y(P_1)=1\)，见结论 \ref{con:xi-terminal-zm-elliptic-rational-points-denominator-law}），并取 \(z_1\in\CC/\Lambda\) 使
\(
P_1=(\wp_\Lambda(z_1),\wp_\Lambda'(z_1)/2)
\)。
由于 \(P_1\) 不在覆盖 \(\pi_y:E\to\mathbb P^1_y\) 的分歧点上（结论 \ref{con:xi-terminal-zm-elliptic-y-hurwitz-passport}），
\(y(z)\) 在 \(z=z_1\) 的邻域内局部可逆；令 \(z=z(y)\) 表示其局部反函数，则主特征支可写为
\[
\lambda_+(y)=X\bigl(z(y)\bigr)=\wp_\Lambda\bigl(z(y)\bigr),
\qquad \lambda_+(1)=2.
\]
因此 \(\lambda_+(y)\) 在避开分歧值集合的任一路径上均可解析延拓，并可通过在 \(\CC/\Lambda\) 上延拓 \(z(y)\) 后回代 \(\wp_\Lambda\) 进行全局表示。
从而 \(\lambda_+(y)\)（以及与之伴随的 \(\log\lambda_+(y)\) 的多值解析结构）可被等价视为固定椭圆格上的反演与周期问题，而非一般代数函数分支的偶然叠加。
进一步，由结论 \ref{con:xi-terminal-zm-leyang-ramification-critical-values} 的微分恒等式
\(
\dd y=(4XY+3X^2-1)\,\omega
\)
（其中 \(\omega=\dd X/(2Y)\)）得：在任意去分歧区域内取满足 \(\pi_y(P(y))=y\) 的局部解析截面 \(P(y)\in E(\CC)\)，令
\(
X(y)=X(P(y))
\)、\(
Y(y)=Y(P(y))
\)，则
\[
\boxed{\;
\frac{\dd}{\dd y}\log \lambda_+(y)
=\frac{1}{X(y)}\frac{\dd X(y)}{\dd y}
=\frac{2Y(y)}{X(y)\,\bigl(4X(y)Y(y)+3X(y)^2-1\bigr)}\; }.
\]
等价地，在最小整模型 \(E_0:\eta^2+\eta=x^3-x\) 上令 \(x=\lambda\)、\(\eta=y-\lambda^2\)，则同一恒等式可写为
\[
\frac{\dd}{\dd y}\log \lambda(y)
=\frac{2\eta+1}{x\bigl(2x(2\eta+1)+3x^2-1\bigr)}.
\]
因此 \(\bigl(\dd(\log\lambda_+)/\dd y\bigr)\dd y=\dd X/X\) 为椭圆曲线上的第三类阿贝尔微分，其周期决定 \(\log\lambda_+\) 的多值解析延拓。
结合注记 \ref{rem:xi-terminal-elliptic-37a1-modular-anchor} 中 \(E\) 作为 \(X_0^+(37)\) 模型的外部锚定，上述周期亦可在 level \(37\) 的模形式周期口径下理解。
其分歧集合与局部型由 \(\pi_y\) 的分歧结构完全确定：有限分歧值为
\(
y\in\{0,1\}\cup\{P_{\mathrm{LY}}(y)=0\}
\)，
且在每个有限分歧值处呈平方根型 Puiseux 分歧（结论 \ref{con:xi-terminal-zm-branch-puiseux-uniformization}）；
而在 \(y=\infty\) 处由 \((y)_\infty=4O\) 给出四重全分歧（结论 \ref{con:xi-terminal-zm-elliptic-y-hurwitz-passport}）。
特别地，实轴上的最近分歧值 \(y_{\mathrm{LY}}\) 及其伴随退化特征值 \(\lambda_{\mathrm{LY}}\) 分别由
\(
P_{\mathrm{LY}}(y_{\mathrm{LY}})=0
\)
与
\(
16\lambda_{\mathrm{LY}}^{3}-9\lambda_{\mathrm{LY}}^{2}+1=0
\)
代数刻画（结论 \ref{con:xi-terminal-zm-leyang-lambda-cubic}）。
\end{remark}

\begin{conclusion}[大偏差速率函数的代数参数化与端点对数奇性]\label{con:xi-terminal-zm-ldp-rate-algebraic-param-endpoint-log}
沿用结论 \ref{con:xi-terminal-zm-ldp} 的压力函数
\(
\psi(s)=\log\lambda_{+}(e^{s})-\log2
\)
及其 Legendre--Fenchel 对偶 \(I(\alpha)=\sup_{s\in\RR}(s\alpha-\psi(s))\)。
对任意 \(y>0\)，令 \(\lambda=\lambda_{+}(y)\)。
定义
\[
\alpha(y):=\frac{y\,\lambda_{+}'(y)}{\lambda_{+}(y)}\in(0,\tfrac12).
\]
则由谱方程 \(\Pi(\lambda,y)=0\) 的隐式微分得到显式有理表达式
\[
\alpha(y)
=-\frac{y\,\partial_y\Pi(\lambda,y)}{\lambda\,\partial_{\lambda}\Pi(\lambda,y)}
=\frac{y\bigl(2\lambda^{2}-(2y+1)\bigr)}{\lambda\bigl(4\lambda^{3}-3\lambda^{2}-2(2y+1)\lambda+1\bigr)}
\Big|_{\lambda=\lambda_{+}(y)}.
\]
并且沿该参数支有
\[
I\bigl(\alpha(y)\bigr)=\alpha(y)\log y-\log\lambda_{+}(y)+\log2.
\]
在本体系中，\(\alpha:(0,\infty)\to(0,\tfrac12)\) 为严格递增的实解析双射，因此 \(I\) 在 \((0,\tfrac12)\) 上实解析且严格凸。
其端点行为由 Perron 分支的分数幂尺度锁定（结论 \ref{con:xi-terminal-zm-perron-puiseux-two-scale}）：
\begin{itemize}
  \item 当 \(\alpha\downarrow0\) 时，
  \[
  I(\alpha)
  =\log2+2\alpha\log\alpha+(\log8-2)\alpha+O\!\bigl(\alpha^{2}\log(1/\alpha)\bigr).
  \]
  \item 当 \(\alpha\uparrow\tfrac12\) 时，令 \(\delta:=\tfrac12-\alpha\downarrow0\)，则
  \[
  I\!\left(\tfrac12-\delta\right)
  =\log2+4\delta\log\delta+(12\log2-4)\delta+O\!\bigl(\delta^{2}\log(1/\delta)\bigr).
  \]
\end{itemize}
其中 \(\alpha\log\alpha\) 与 \(\delta\log\delta\) 两个非解析项分别由 \(y\to0^{+}\) 的 \(y^{1/2}\) Puiseux 与 \(y\to+\infty\) 的 \(y^{-1/4}\) Puiseux 诱导，因而构成该 Lee--Yang 权重体系在两端点的可检验奇性指纹。
\end{conclusion}

\begin{conclusion}[复 Lee--Yang 主导的解析半径与高阶累积量振荡律]\label{con:xi-terminal-zm-psi-analytic-radius-cumulant-oscillation}
考虑 \(\psi(s)=\log\lambda_{+}(e^{s})-\log2\) 在 \(s=0\) 的解析芽（结论 \ref{con:xi-terminal-zm-lln-clt}）。
由结论 \ref{con:xi-terminal-zm-discriminant-leyang} 与结论 \ref{con:xi-terminal-zm-leyang-cubic-branch-image}，
\(y\)-平面的有限分歧值除 \(\{0,1\}\) 外还包含 Lee--Yang 三次
\(
P_{\mathrm{LY}}(y)=256y^{3}+411y^{2}+165y+32
\)
的共轭复根对
\[
y_{\mathrm c}^{\pm}\approx-0.2355083735\pm0.2339255521\,i.
\]
令
\[
s_{\mathrm c}^{\pm}:=\log y_{\mathrm c}^{\pm}\quad(\text{取主值}),
\qquad
\rho:=|s_{\mathrm c}^{+}|\approx2.6045558274,
\qquad
\theta:=\arg(s_{\mathrm c}^{+})\approx2.0080026785,
\]
则 \(\psi\) 在圆盘 \(|s|<\rho\) 内单值全纯，并且该圆盘为以 \(0\) 为中心的最大解析圆盘；其边界上首先遇到的奇点即 \(s_{\mathrm c}^{\pm}\) 所对应的代数分歧。
\par
记 \(\kappa_n^\infty:=\psi^{(n)}(0)\in\QQ\) 为单位长度的第 \(n\) 阶累积量密度（结论 \ref{con:xi-terminal-zm-cumulant-rational}）。
则当 \(n\to\infty\) 时存在常数 \(B\neq 0\) 与相位 \(\varphi\in\RR\)，使得
\[
\kappa_n^\infty
=\frac{n!}{\sqrt\pi}\,|B|\,\rho^{-n+\frac12}\,n^{-\frac32}
\cos\!\Bigl((n-\tfrac12)\theta+\varphi\Bigr)
+o\!\bigl(n!\rho^{-n}n^{-\frac32}\bigr).
\]
更精确地，取临界对 \((\lambda_{\mathrm c},y_{\mathrm c})\) 满足
\[
\Pi(\lambda_{\mathrm c},y_{\mathrm c})=0,\qquad
\partial_{\lambda}\Pi(\lambda_{\mathrm c},y_{\mathrm c})=0,
\qquad (y_{\mathrm c}=y_{\mathrm c}^{+}),
\]
并令
\[
c:=\sqrt{-\,\frac{2\,\partial_{y}\Pi(\lambda_{\mathrm c},y_{\mathrm c})}{\partial_{\lambda\lambda}\Pi(\lambda_{\mathrm c},y_{\mathrm c})}},
\qquad
B:=\frac{c\sqrt{y_{\mathrm c}}}{\lambda_{\mathrm c}},
\]
则上述振荡律成立；数值上可取
\[
\lambda_{\mathrm c}\approx0.4179674053-0.2321140339\,i,
\qquad
|B|\approx0.9634999692.
\]
因此高阶累积量密度除呈现 \(\rho^{-n}\) 的指数控制与 \(n^{-3/2}\) 的普适幂律修正外，还携带由 \(\theta=\arg(s_{\mathrm c}^{+})\) 决定的确定性振荡频率，给出一条可直接数值核验的“复 Lee--Yang 振荡律”。
可复现核验脚本：\texttt{scripts/exp\_fold\_zm\_bivariate\_partition\_audit.py}。
\end{conclusion}

\begin{theorem}[Lee--Yang 主振荡角的 Gelfond--Schneider 超越性]\label{thm:xi-terminal-zm-leyang-oscillation-angle-transcendence}
沿用结论 \ref{con:xi-terminal-zm-psi-analytic-radius-cumulant-oscillation} 的记号，令
\(
s_{\mathrm c}:=s_{\mathrm c}^{+}=\log y_{\mathrm c}^{+}
\)
为主值对数，并写 \(\rho:=|s_{\mathrm c}|\)、\(\theta:=\arg(s_{\mathrm c})\in(0,\pi)\)。
则：
\begin{enumerate}
  \item \(\tan\theta\) 为超越数；
  \item 特别地，\(\theta/\pi\notin\QQ\)。
\end{enumerate}
\end{theorem}
\begin{proof}
设
\[
N:=y_{\mathrm c}^{+}\overline{y_{\mathrm c}^{+}}=|y_{\mathrm c}^{+}|^{2}\in\overline{\QQ}\cap\RR_{>0},
\qquad
r:=\frac{y_{\mathrm c}^{+}}{\overline{y_{\mathrm c}^{+}}}\in\overline{\QQ},
\qquad |r|=1.
\]
由于 \(y_{\mathrm c}^{+}\notin\RR_{\le 0}\)，主值对数与共轭相容，故
\[
\log N=s_{\mathrm c}+\overline{s_{\mathrm c}}=2\Re(s_{\mathrm c})\neq 0,
\qquad
s_{\mathrm c}-\overline{s_{\mathrm c}}=2i\,\Im(s_{\mathrm c})
\]
并且 \(\exp(s_{\mathrm c}-\overline{s_{\mathrm c}})=r\neq 1\)。
令
\[
\beta:=\frac{s_{\mathrm c}-\overline{s_{\mathrm c}}}{s_{\mathrm c}+\overline{s_{\mathrm c}}}
=i\,\frac{\Im(s_{\mathrm c})}{\Re(s_{\mathrm c})}=i\tan\theta.
\]
若 \(\tan\theta\) 为代数数，则 \(\beta\in\overline{\QQ}\setminus\QQ\)（因为 \(\beta\notin\RR\)）。
取 \(\alpha:=N\in\overline{\QQ}\setminus\{0,1\}\)，由 Gelfond--Schneider 定理可知任一值 \(\alpha^{\beta}\) 皆为超越数。\cite{Baker1975TranscendentalNumberTheory}
另一方面，由 \(\beta\log N=s_{\mathrm c}-\overline{s_{\mathrm c}}\) 得
\[
N^{\beta}=\exp(\beta\log N)=\exp(s_{\mathrm c}-\overline{s_{\mathrm c}})=r\in\overline{\QQ},
\]
矛盾。
故 \(\tan\theta\) 不可能为代数数，即 \(\tan\theta\) 为超越数。
\par
若 \(\theta/\pi\in\QQ\)，则 \(\zeta:=e^{2i\theta}\) 为单位根，从而
\(
\tan\theta=\frac{\zeta-1}{i(\zeta+1)}
\in\overline{\QQ}
\)，
与已证超越性矛盾。
故 \(\theta/\pi\notin\QQ\)。证毕。
\end{proof}

\begin{theorem}[归一化高阶累积量密度的反正弦极限律]\label{thm:xi-terminal-zm-cumulant-arcsine-limit-law}
\leanverified{paper\_xi\_terminal\_zm\_cumulant\_arcsine\_limit\_law}
沿用结论 \ref{con:xi-terminal-zm-psi-analytic-radius-cumulant-oscillation} 的振荡渐近：
存在 \(B\neq 0\)、\(\varphi\in\RR\) 使
\[
\kappa_n^\infty
=\frac{n!}{\sqrt\pi}\,|B|\,\rho^{-n+\frac12}\,n^{-\frac32}
\cos\!\Bigl((n-\tfrac12)\theta+\varphi\Bigr)
+o\!\bigl(n!\rho^{-n}n^{-\frac32}\bigr).
\]
定义归一化序列
\[
A_n:=\frac{\sqrt\pi}{|B|}\,\rho^{\,n-\frac12}\,n^{\frac32}\,\frac{\kappa_n^\infty}{n!}.
\]
则 \(A_n=\cos\!\bigl((n-\tfrac12)\theta+\varphi\bigr)+o(1)\)，并且由定理 \ref{thm:xi-terminal-zm-leyang-oscillation-angle-transcendence} 的 \(\theta/\pi\notin\QQ\) 得：\(A_n\) 的经验分布弱收敛到 \([-1,1]\) 上的反正弦分布。
具体地，对任意连续函数 \(f\in C([-1,1])\)，有
\[
\lim_{N\to\infty}\frac1N\sum_{n=1}^{N} f(A_n)
=\frac{1}{\pi}\int_{-1}^{1}\frac{f(x)}{\sqrt{1-x^2}}\,\dd x.
\]
\end{theorem}
\begin{proof}
由渐近式直接得到 \(A_n-\cos(t_n)\to 0\)，其中 \(t_n:=(n-\tfrac12)\theta+\varphi\)。
由 \(\theta/\pi\notin\QQ\) 可知旋转序列 \(t_n\ (\mathrm{mod}\ 2\pi)\) 在 \(\RR/2\pi\ZZ\) 上等分布（Kronecker--Weyl），故对任意连续 \(g\in C(\RR/2\pi\ZZ)\) 有
\[
\frac1N\sum_{n=1}^{N} g(t_n)\to \frac{1}{2\pi}\int_{0}^{2\pi} g(t)\,\dd t.
\]
取 \(g(t)=f(\cos t)\) 即得 \(\cos(t_n)\) 的经验分布为反正弦律的推前测度。\cite{KuipersNiederreiter1974}
再由 \(f\) 的一致连续性与 \(A_n-\cos(t_n)\to 0\) 控制 Ces\`aro 平均差，得到所述极限式。证毕。
\end{proof}

\begin{corollary}[符号密度、极值充满与无限次符号翻转]\label{cor:xi-terminal-zm-cumulant-sign-density-extremes}
\leanverified{paper\_xi\_terminal\_zm\_cumulant\_sign\_density\_extremes}
沿用定理 \ref{thm:xi-terminal-zm-cumulant-arcsine-limit-law} 的记号。
则集合 \(\{n:\kappa_n^\infty>0\}\) 与 \(\{n:\kappa_n^\infty<0\}\) 的自然密度均为 \(1/2\)，并且
\[
\limsup_{n\to\infty}A_n=1,\qquad \liminf_{n\to\infty}A_n=-1.
\]
从而序列 \((\kappa_n^\infty)\) 发生无限次符号翻转，且不可能最终周期。
\end{corollary}
\begin{proof}
由反正弦分布关于 \(0\) 对称且 \(\cos t\) 在一半圆周取正、一半取负，
结合 \(A_n-\cos(t_n)\to 0\) 得符号密度结论。
反正弦分布的支撑为 \([-1,1]\) 且任意开邻域均有正测度，故 \(\cos(t_n)\) 在 \([-1,1]\) 稠密；
再由 \(A_n-\cos(t_n)\to 0\) 得 \(\limsup A_n=1\)、\(\liminf A_n=-1\)。
最后，无限次正负出现推出无限次符号翻转；若 \((\kappa_n^\infty)\) 最终周期，则其主振荡相位 \((n-\tfrac12)\theta+\varphi\) 只能对应有理旋转，从而 \(\theta/\pi\in\QQ\)，与定理 \ref{thm:xi-terminal-zm-leyang-oscillation-angle-transcendence} 矛盾。证毕。
\end{proof}

\begin{proposition}[相邻阶局部增长率在扩展实轴上的稠密聚点]\label{prop:xi-terminal-zm-cumulant-local-growth-dense}
\leanverified{paper\_xi\_terminal\_zm\_cumulant\_local\_growth\_dense}
沿用结论 \ref{con:xi-terminal-zm-psi-analytic-radius-cumulant-oscillation} 的振荡渐近，并在 \(\kappa_n^\infty\neq 0\) 时定义
\[
R_n:=\rho\left(\frac{n}{n+1}\right)^{\frac32}\cdot\frac{\kappa_{n+1}^\infty}{(n+1)\kappa_n^\infty}.
\]
则 \(R_n\) 在扩展实轴 \(\widehat{\RR}:=\RR\cup\{\infty\}\) 上具有稠密的聚点集合：对任意 \(L\in\RR\)，存在无穷子列 \(n_k\to\infty\) 使 \(R_{n_k}\to L\)，并且亦存在无穷子列使 \(|R_{n_k}|\to\infty\)。
\end{proposition}
\begin{proof}
由结论 \ref{con:xi-terminal-zm-psi-analytic-radius-cumulant-oscillation}，
\[
\frac{\kappa_{n+1}^\infty}{(n+1)\kappa_n^\infty}
=\frac{1}{\rho}\left(\frac{n}{n+1}\right)^{\frac32}
\frac{\cos(t_n+\theta)}{\cos(t_n)}+o(1),
\qquad
t_n:=(n-\tfrac12)\theta+\varphi.
\]
故
\[
R_n=\frac{\cos(t_n+\theta)}{\cos(t_n)}+o(1)
=\cos\theta-\sin\theta\,\tan(t_n)+o(1),
\]
其中 \(\sin\theta\neq 0\)（因 \(\theta\in(0,\pi)\)）。
由 \(\theta/\pi\notin\QQ\) 得 \(t_n\ (\mathrm{mod}\ \pi)\) 等分布，从而 \(\tan(t_n)\) 在 \(\widehat{\RR}\) 上稠密；
再由仿射同胚 \(x\mapsto \cos\theta-\sin\theta\,x\) 与 \(o(1)\) 微扰的稳定性得到 \(R_n\) 的聚点稠密性，并可取 \(|\tan(t_n)|\to\infty\) 的子列推出 \(|R_n|\to\infty\)。证毕。
\end{proof}

\begin{theorem}[Stokes 幅值平方 \texorpdfstring{$B^2$}{B^2} 的三次最小多项式、判别式与 \texorpdfstring{$S_3$}{S3} 刚性]\label{thm:xi-terminal-zm-stokes-amplitude-square-cubic-rigidity}
沿用结论 \ref{con:xi-terminal-zm-psi-analytic-radius-cumulant-oscillation} 的临界对 \((\lambda_{\mathrm c},y_{\mathrm c})\) 与幅值常数 \(B\)，并令
\[
b:=B^{2}=\frac{c^{2}y_{\mathrm c}}{\lambda_{\mathrm c}^{2}}\in\CC,
\qquad t:=b.
\]
设 \(K:=\QQ(y_{\mathrm c})\) 为 Lee--Yang 三次域（\(P_{\mathrm{LY}}(y_{\mathrm c})=0\)）。
则在 \(K\) 内成立显式有理恒等式
\begin{equation}\label{eq:xi-terminal-zm-stokes-t-rational-in-yc}
\boxed{\;
t
=-\frac{1}{6}\cdot
\frac{440y_{\mathrm c}^{2}+47y_{\mathrm c}-64}{46y_{\mathrm c}^{2}+40y_{\mathrm c}-5}\;}
\end{equation}
并且 \(t\) 在 \(\QQ\) 上满足三次整式
\begin{equation}\label{eq:xi-terminal-zm-stokes-t-minpoly}
\boxed{\;
162t^{3}+1593t^{2}+1998t+1184=0\;}
\end{equation}
等价地，\(b\) 满足
\[
162b^{3}+1593b^{2}+1998b+1184=0.
\]
该三次在 \(\QQ\) 上不可约，且判别式为
\[
\mathrm{Disc}(162b^{3}+1593b^{2}+1998b+1184)
=-2^{2}\cdot 3^{9}\cdot 11^{2}\cdot 37\cdot 109^{2}.
\]
因判别式非平方，其分裂域伽罗瓦群同构于 \(S_3\)。
从而 \(\QQ(b)=\QQ(t)=\QQ(y_{\mathrm c})\)，亦即 \(B^2\) 本身为 \(K\) 的一个原始生成元。
\end{theorem}
\begin{proof}[证明要点]
由结论 \ref{con:xi-terminal-zm-branch-double-root-linear-closure}，在约束 \(P_{\mathrm{LY}}(y_{\mathrm c})=0\) 下二重根谱支满足线性刚性
\(
279\lambda_{\mathrm c}+512y_{\mathrm c}^{2}+518y_{\mathrm c}+5=0
\)，
故 \(\lambda_{\mathrm c}\in \QQ(y_{\mathrm c})\)。
另一方面，
\(
\partial_y\Pi(\lambda,y)=-2\lambda^2+2y+1
\)、
\(
\partial_{\lambda\lambda}\Pi(\lambda,y)=12\lambda^2-6\lambda-4y-2
\)，
从而
\(
t=-2\cdot\frac{y_{\mathrm c}\,\partial_y\Pi(\lambda_{\mathrm c},y_{\mathrm c})}{\lambda_{\mathrm c}^{2}\,\partial_{\lambda\lambda}\Pi(\lambda_{\mathrm c},y_{\mathrm c})}
\in \QQ(\lambda_{\mathrm c},y_{\mathrm c})
=\QQ(y_{\mathrm c})
\)。
将 \(\lambda_{\mathrm c}\) 以 \(y_{\mathrm c}\) 的有理式表示并在三次商环 \(\QQ[y]/(P_{\mathrm{LY}})\) 中约化，即得 \eqref{eq:xi-terminal-zm-stokes-t-rational-in-yc}。
再将 \eqref{eq:xi-terminal-zm-stokes-t-rational-in-yc} 改写为关于 \(y_{\mathrm c}\) 的二次方程并与 \(P_{\mathrm{LY}}(y_{\mathrm c})=0\) 联立作消元，即得到 \eqref{eq:xi-terminal-zm-stokes-t-minpoly}。
有理根检验排除一次因子，故该三次在 \(\QQ\) 上不可约。
其判别式按三次判别式公式计算即为
\(
-2^{2}\cdot 3^{9}\cdot 11^{2}\cdot 37\cdot 109^{2}
\)，非平方，故分裂域伽罗瓦群为 \(S_3\)。
由不可约性可知 \([\QQ(t):\QQ]=3\)，且 \(t\in\QQ(y_{\mathrm c})\) 已证，遂 \(\QQ(t)=\QQ(y_{\mathrm c})\)；
由 \(b=t\) 得 \(\QQ(b)=\QQ(t)\)。证毕。
\end{proof}

\begin{corollary}[幅值三次的二次分辨子域：\texorpdfstring{$\QQ(\sqrt{-111})$}{Q(sqrt(-111))}]\label{cor:xi-terminal-zm-stokes-b2-quadratic-resolvent}
沿用定理 \ref{thm:xi-terminal-zm-stokes-amplitude-square-cubic-rigidity} 的记号，令 \(K=\QQ(b)\) 并令 \(L\) 为 \(K/\QQ\) 的伽罗瓦闭包。
则 \(L/\QQ\) 的唯一二次子域为
\[
\boxed{\ \QQ(\sqrt{-111})\ }.
\]
特别地，该二次子域与 Lee--Yang 三次的二次分辨域一致（见结论 \ref{con:xi-terminal-zm-leyang-cubic-arithmetic}），从而 \(B^2\) 与 \(y_{\mathrm c}\) 的任一原始生成元给出同一条二次判别脊柱。
\end{corollary}
\begin{proof}
由定理 \ref{thm:xi-terminal-zm-stokes-amplitude-square-cubic-rigidity}，\(b\) 的极小多项式为不可约三次 \(f(t)=162t^{3}+1593t^{2}+1998t+1184\)，且
\[
\mathrm{Disc}(f)=-2^{2}\cdot 3^{9}\cdot 11^{2}\cdot 37\cdot 109^{2}.
\]
判别式的平方自由部分为 \(-3\cdot 37=-111\)。
对不可约三次且 \(\mathrm{Disc}(f)\) 非平方的情形，其分裂域的伽罗瓦群为 \(S_3\)，并且其唯一二次子域由 \(\QQ(\sqrt{\mathrm{Disc}(f)})\) 给出；
因此 \(L\) 的唯一二次子域为 \(\QQ(\sqrt{-111})\)。证毕。
\end{proof}

\begin{conjecture}[Stokes 幅值平方的 CM--模参数化]\label{conj:xi-terminal-zm-stokes-b2-cm-modular-parameterization}
沿用定理 \ref{thm:xi-terminal-zm-stokes-amplitude-square-cubic-rigidity} 的记号，令 \(b:=B^2\) 并记三次域 \(K:=\QQ(b)\)。
猜想存在某一层级 \(N\) 的模曲线上的模函数 \(f\)（例如 \(X_0(N)\) 的 Hauptmodul 或 \(\lambda\)-函数的有理变换）与某一 CM 点 \(\tau\in\mathbb H\)，使得
\[
b=f(\tau),
\qquad
\QQ(b)\ \text{为相应环类域的某一三次子域}.
\]
并且 \(N\) 的素因子集合可由 \(\mathrm{Disc}(162b^{3}+1593b^{2}+1998b+1184)\) 的素因子集合控制。
该猜想可通过对候选 \(N\) 与 CM 判别式的有限枚举检验：比较 \(f(\tau)\) 的极小多项式与 \eqref{eq:xi-terminal-zm-stokes-t-minpoly}。
\end{conjecture}

% Exponential period model for the Stokes amplitude via third-kind differentials.

\begin{conjecture}[Stokes 幅值平方的第三类微分指数化周期模型]\label{conj:xi-terminal-zm-stokes-b2-thirdkind-exponential-period}
沿用定理 \ref{thm:xi-terminal-zm-leyang-elliptic-structure} 的椭圆曲线 \(E_{\min}/\QQ\) 及其有理函数
\(
y= X^{2}+Y-\tfrac12
\)
（在短 Weierstrass 口径 \(Y^{2}=X^{3}-X+\tfrac14\) 下），并沿用定理 \ref{thm:xi-terminal-zm-stokes-amplitude-square-cubic-rigidity} 的幅值常数 \(B\)。
则存在一个零度因子 \(D\in\mathrm{Div}^{0}(E_{\min})\)，其支撑包含于
\[
\mathrm{supp}(D)\subseteq y^{-1}(0)\cup y^{-1}(1),
\]
以及一条由 Stokes 归一化选择所确定的基本闭路 \(\gamma\in H_{1}(E_{\min}(\CC),\ZZ)\)，使得对满足
\(
\mathrm{Res}(\eta_D)=D
\)
且 \(a\)-周期为零的第三类微分 \(\eta_D\)，成立指数化周期恒等式
\[
\boxed{\;
B^{2}=\exp\!\left(\int_{\gamma}\eta_D\right).\;}
\]
进一步地，若 \(\varphi:X_0(37)\to E_{\min}\) 为模参数化（同构到 \(37\mathrm a1\) 的模性锚点口径），并令 \(f\in S_2(\Gamma_0(37))\) 为与 \(E_{\min}\) 对应的权二新形式，则预期 \(\varphi^{*}\eta_D\) 属于由 \(f(q)\,\dd q/q\) 及其 \(q\)-导数所张成的微分域（等价地，可由该模形式及其微分闭包中的元素表示）。
\end{conjecture}

\endinput



\begin{theorem}[振幅平方 \texorpdfstring{$B^2$}{B^2} 与 \texorpdfstring{$\kappa_*^2$}{kappa*^2} 的显式有理互构]\label{thm:xi-terminal-zm-b2-kappa2-rational-interconstruction}
\leanverified{paper\_xi\_terminal\_zm\_b2\_kappa2\_rational\_interconstruction}
令
\[
u:=\kappa_*^2,
\qquad
b:=B^2.
\]
其中 \(u\) 取自定理 \ref{thm:xi-terminal-zm-kappa-square-cubic-field-s3}，\(b\) 取自定理 \ref{thm:xi-terminal-zm-stokes-amplitude-square-cubic-rigidity}。
则存在显式有理关系
\[
\boxed{\;
b
=-2\cdot
\frac{24708u^2-28628u+7733}{51084u^2-44412u+8735}\; }.
\]
因而
\[
\QQ(b)=\QQ(u)=\QQ(\lambda_*),
\]
即 \(B^2\) 与 \(\kappa_*^2\) 生成同一 Lee--Yang 三次数域。
\end{theorem}
\begin{proof}
由定理 \ref{thm:xi-terminal-zm-kappa-square-cubic-field-s3}，
\[
\lambda_*=\frac{3(2u-1)}{4(3u-2)}.
\]
由定理 \ref{thm:xi-terminal-zm-stokes-t-elliptic-branch-coordinate} 的表达式 \eqref{eq:xi-terminal-zm-stokes-t-rational-in-alpha}（取 \(\alpha=\lambda_*\) 且 \(t=b\)）：
\[
b=\frac{2(1113\lambda_*^2-2192\lambda_*+303)}{3(1831\lambda_*^2+144\lambda_*-175)}.
\]
把 \(\lambda_*=\frac{3(2u-1)}{4(3u-2)}\) 代入并在 \(\QQ(u)\) 中约化，即得所示有理式，故 \(b\in\QQ(u)\)。
另一方面，定理 \ref{thm:xi-terminal-zm-kappa-square-cubic-field-s3} 与定理 \ref{thm:xi-terminal-zm-stokes-amplitude-square-cubic-rigidity} 分别给出
\(
[\QQ(u):\QQ]=[\QQ(b):\QQ]=3
\)。
故包含关系 \(\QQ(b)\subseteq\QQ(u)\) 只能是等号；再结合
\(
\QQ(u)=\QQ(\lambda_*)
\)
即得结论。
\end{proof}

\begin{theorem}[幅值平方的椭圆分歧坐标表达]\label{thm:xi-terminal-zm-stokes-t-elliptic-branch-coordinate}
沿用定理 \ref{thm:xi-terminal-zm-stokes-amplitude-square-cubic-rigidity} 的记号，并令 \(\alpha:=\lambda_{\mathrm c}\)。
则 \(\alpha\) 满足 Lee--Yang 三次谱分歧方程 \(16\alpha^{3}-9\alpha^{2}+1=0\)（结论 \ref{con:xi-terminal-zm-leyang-ramification-critical-values}）。
并且在分歧约束下有恒等式
\[
4\alpha\,y_{\mathrm c}=4\alpha^{3}-3\alpha^{2}-2\alpha+1
\qquad(\text{结论 \ref{con:xi-terminal-zm-leyang-ramification-critical-values}}),
\]
从而 \(t=B^{2}\) 可写为 \(\QQ(\alpha)\) 中的有理函数：
\begin{equation}\label{eq:xi-terminal-zm-stokes-t-rational-in-alpha}
\boxed{\;
t
=\frac{2\,(1113\alpha^{2}-2192\alpha+303)}{3\,(1831\alpha^{2}+144\alpha-175)}\; }.
\end{equation}
\end{theorem}
\begin{proof}[证明要点]
由结论 \ref{con:xi-terminal-zm-leyang-ramification-critical-values} 的分歧恒等式可将 \(y_{\mathrm c}\) 表为 \(\alpha\) 的有理式。
将其代入定理 \ref{thm:xi-terminal-zm-stokes-amplitude-square-cubic-rigidity} 的表达式 \eqref{eq:xi-terminal-zm-stokes-t-rational-in-yc}，
并在商环 \(\QQ[\alpha]/(16\alpha^{3}-9\alpha^{2}+1)\) 中约化，即得 \eqref{eq:xi-terminal-zm-stokes-t-rational-in-alpha}。证毕。
\end{proof}

\begin{theorem}[隐藏平方与纯有理反演：由 \texorpdfstring{$t=B^{2}$}{t=B^2} 重构 \texorpdfstring{$y_{\mathrm c}$}{yc}]\label{thm:xi-terminal-zm-stokes-hidden-square-rational-inversion}
令 \(t\) 满足 \eqref{eq:xi-terminal-zm-stokes-t-minpoly}，并在三次域 \(K:=\QQ(t)\) 中定义
\begin{equation}\label{eq:xi-terminal-zm-stokes-s-def}
s:=\frac{-71045-13947t+12258t^{2}}{1199}.
\end{equation}
则在 \(K\) 内成立隐藏平方恒等式
\begin{equation}\label{eq:xi-terminal-zm-stokes-hidden-square}
\boxed{\;
s^{2}=10080t^{2}+16224t+12761\; }.
\end{equation}
并且由 \eqref{eq:xi-terminal-zm-stokes-t-rational-in-yc} 导出的二次方程在 \(K\) 内分解，其两根可写为纯有理表达
\begin{equation}\label{eq:xi-terminal-zm-stokes-yc-rational-in-t}
\boxed{\;
y^{(1)}(t)
=\frac{36774t^{2}-329601t-269488}{9592(69t+110)},
\qquad
y^{(2)}(t)
=\frac{-36774t^{2}-245919t+156782}{9592(69t+110)}\; }.
\end{equation}
其中与 \(P_{\mathrm{LY}}(y)=0\) 相容者即为 Lee--Yang 临界值 \(y_{\mathrm c}\)；在结论 \ref{con:xi-terminal-zm-psi-analytic-radius-cumulant-oscillation} 取 \(y_{\mathrm c}=y_{\mathrm c}^{+}\) 的嵌入下，对应 \(y_{\mathrm c}=y^{(1)}(t)\)。
\end{theorem}
\begin{proof}[证明要点]
由 \eqref{eq:xi-terminal-zm-stokes-t-rational-in-yc} 可得 \(y_{\mathrm c}\) 满足关于 \(y\) 的二次方程
\[
(276t+440)y^{2}+(240t+47)y-(30t+64)=0.
\]
其判别式为
\(
\Delta=9(10080t^{2}+16224t+12761)
\)。
式 \eqref{eq:xi-terminal-zm-stokes-hidden-square} 表明该判别式平方根已属于 \(K\)，故该二次方程在 \(K\) 内分解并给出两条根的纯有理表达；
化简即得 \eqref{eq:xi-terminal-zm-stokes-yc-rational-in-t}。证毕。
\end{proof}

\begin{proposition}[临界谱值的幅值重构]\label{prop:xi-terminal-zm-stokes-lambda-critical-rational-in-t}
在定理 \ref{thm:xi-terminal-zm-stokes-hidden-square-rational-inversion} 的设定下，
\(\lambda_{\mathrm c}\) 可由 \(t\) 显式重构为
\begin{equation}\label{eq:xi-terminal-zm-stokes-lambda-critical-rational-in-t}
\boxed{\;
\lambda_{\mathrm c}
=-\frac{525411t^{2}-4582128t-2625484}{1199(69t+110)^{2}}\; }.
\end{equation}
\end{proposition}
\begin{proof}
由结论 \ref{con:xi-terminal-zm-branch-double-root-linear-closure} 的线性刚性关系，
\(
\lambda_{\mathrm c}=-(512y_{\mathrm c}^{2}+518y_{\mathrm c}+5)/279
\)。
将 \eqref{eq:xi-terminal-zm-stokes-yc-rational-in-t} 中与 \(P_{\mathrm{LY}}(y)=0\) 相容的根 \(y_{\mathrm c}\) 代入并在 \(\QQ(t)\) 中约化即可得到 \eqref{eq:xi-terminal-zm-stokes-lambda-critical-rational-in-t}。证毕。
\end{proof}

\begin{theorem}[Stokes 幅值的六次代数度与极小多项式]\label{thm:xi-terminal-zm-stokes-amplitude-degree6-minpoly}
设 \(B\) 为结论 \ref{con:xi-terminal-zm-psi-analytic-radius-cumulant-oscillation} 中由临界对 \((\lambda_{\mathrm c},y_{\mathrm c})\) 所定义的 Stokes 幅值常数，并令 \(t=B^{2}\)。
则 \([\QQ(B):\QQ]=6\)，且 \(B\) 的极小多项式为
\begin{equation}\label{eq:xi-terminal-zm-stokes-B-minpoly}
\boxed{\;
162B^{6}+1593B^{4}+1998B^{2}+1184=0\; }.
\end{equation}
\end{theorem}
\begin{proof}
由 \eqref{eq:xi-terminal-zm-stokes-t-minpoly}，\(t\) 在 \(\QQ\) 上次数为 \(3\)，故 \(K:=\QQ(t)\) 为三次域且 \(K\subset \QQ(B)\)。
若 \(B\in K\)，则 \(t=B^{2}\) 在 \(K\) 中为平方，从而其范数 \(\mathrm{N}_{K/\QQ}(t)\) 必为有理数平方，因而非负。
另一方面，把 \eqref{eq:xi-terminal-zm-stokes-t-minpoly} 除以 \(162\) 化为首一多项式，可得
\[
\mathrm{N}_{K/\QQ}(t)=(-1)^{3}\cdot\frac{1184}{162}=-\frac{592}{81}<0,
\]
矛盾。
故 \(B\notin K\)，从而 \([\QQ(B):K]=2\) 并推出 \([\QQ(B):\QQ]=6\)。
式 \eqref{eq:xi-terminal-zm-stokes-B-minpoly} 由 \eqref{eq:xi-terminal-zm-stokes-t-minpoly} 代入 \(t=B^{2}\) 直接得到；
其次数为 \(6\)，与 \([\QQ(B):\QQ]=6\) 一致，故为极小多项式。证毕。
\end{proof}

\begin{conclusion}[全阶累积量密度的有理性与可审计生成方程]\label{con:xi-terminal-zm-cumulant-dalgebraic-equation}
结论 \ref{con:xi-terminal-zm-cumulant-rational} 已给出 \(\psi(s)\) 在 \(s=0\) 的 Taylor 展开属于 \(\QQ[[s]]\)，因而 \(\kappa_n^\infty=\psi^{(n)}(0)\in\QQ\) 对一切 \(n\ge 1\) 成立。
进一步地，由 \(\Pi(\lambda,y)=0\) 在代换
\(
\lambda=2e^{\psi(s)},\ y=e^{s}
\)
下得到 \(\psi\) 的显式代数生成方程
\[
16e^{4\psi(s)}-8e^{3\psi(s)}-4(2e^{s}+1)e^{2\psi(s)}+2e^{\psi(s)}+e^{2s}+e^{s}=0,
\]
从而 \(\psi\) 为 \(D\)-代数函数；这为任意阶 \(\kappa_n^\infty\) 的符号递推与误差审计提供封闭的代数约束。
除结论 \ref{con:xi-terminal-zm-cumulant-rational} 已列出的低阶项外，高阶例如
\[
\kappa_{5}^\infty=\frac{15650}{1594323},\quad
\kappa_{6}^\infty=\frac{307723}{57395628},\quad
\kappa_{7}^\infty=-\frac{2332358}{129140163},\quad
\kappa_{8}^\infty=\frac{166894943}{55788550416}.
\]
可复现核验脚本：\texttt{scripts/exp\_fold\_zm\_bivariate\_partition\_audit.py}。
\end{conclusion}
